% GENERATED FILE. Edit canonical lessons, then run npm run generate:books.
% canonical-source-hash: fnv1a64:187f8fbfacfebab6
% canonical-lessons: GU-C153-ordo, GU-C153-rasodu, GU-C153-agasi, GU-C153-osri, GU-C153-dadar

\chapter{Rooms, roof and veranda}
\label{ch:gu-rooms-roof-and-veranda}
\hlchaptermodality{153}

\begin{chapteropening}
\textbf{By the end of this chapter:} \emph{I can name the rooms and spaces of a Gujarati home.}

Complete the last lesson of chapter 153: I can name the rooms and spaces of a Gujarati home.
\end{chapteropening}

\section[\texorpdfstring{\gu{ઓરડો} (orḍo)}{ઓરડો (orḍo)}]{\gu{ઓરડો} (orḍo) — room}

\label{lesson:GU-C153-ordo}



\begin{quote}
Before the new one: say the Gujarati for to believe, then the Gujarati for to change.
\end{quote}

\subsection*{You'll want to know: \gu{ઓરડો}}
\textbf{\gu{ઓરડો}} — \emph{orḍo} — \textquotedblleft{}room\textquotedblright{}.

\begin{quote}
\textbf{\gu{આ} \gu{મારો} \gu{ઓરડો} \gu{છે}.} — \emph{ā māro orḍo che} — \textquotedblleft{}This is my room.\textquotedblright{}
\end{quote}

\begin{cousinweb}
\textbf{\gu{ઓરડો}} (\emph{orḍo}) is a tadbhava, a word inherited from Sanskrit and reshaped by sound change. It comes from \emph{apavaraka}, \textquotedblleft{}inner apartment\textquotedblright{}, and CDIAL glosses the Gujarati word \textquotedblleft{}principal room\textquotedblright{} (CDIAL 450).
\end{cousinweb}

\begin{grammarlens}[title={What you've built}]
One more word for everyday life at A2.
\end{grammarlens}

\subsection*{Your turn}
\begin{itemize}
\raggedright
  \item \emph{Say it:} \emph{orḍo}
  \item \emph{Say it:} \emph{orḍo}, once more
  \item \emph{Say it:} \emph{ā māro orḍo che} — This is my room.
  \item \emph{Recall:} say the Gujarati for to believe, then the Gujarati for to change, then say \emph{orḍo} again
\end{itemize}

\begin{tcolorbox}[breakable,colback=teal!4,colframe=teal!35!black,title={Before you move on}]
What does \gu{ઓરડો} mean? (\textquotedblleft{}Room\textquotedblright{}.) Say it once more.
\end{tcolorbox}
\section[\texorpdfstring{\gu{રસોડું} (rasoḍũ)}{રસોડું (rasoḍũ)}]{\gu{રસોડું} (rasoḍũ) — kitchen}

\label{lesson:GU-C153-rasodu}



\begin{quote}
Before the new one: say the Gujarati for to change, then the Gujarati for room.
\end{quote}

\subsection*{You'll want to know: \gu{રસોડું}}
\textbf{\gu{રસોડું}} — \emph{rasoḍũ} — \textquotedblleft{}kitchen\textquotedblright{}.

\begin{quote}
\textbf{\gu{માતા} \gu{રસોડામાં} \gu{રાંધે} \gu{છે}.} — \emph{mātā rasoḍāmā̃ rā̃dhe che} — \textquotedblleft{}Mother cooks in the kitchen.\textquotedblright{}
\end{quote}

\begin{cousinweb}
\textbf{\gu{રસોડું}} (\emph{rasoḍũ}) is a tadbhava. It comes from \emph{*rasavāṭa}, \textquotedblleft{}cooking enclosure\textquotedblright{}, a reconstructed form built from \emph{rasa}, \textquotedblleft{}juice, flavour\textquotedblright{}, and \emph{vāṭa}, \textquotedblleft{}enclosure\textquotedblright{}; CDIAL glosses the Gujarati word \textquotedblleft{}kitchen\textquotedblright{} (CDIAL 10658).
\end{cousinweb}

\begin{grammarlens}[title={What you've built}]
One more word for everyday life at A2.
\end{grammarlens}

\subsection*{Your turn}
\begin{itemize}
\raggedright
  \item \emph{Say it:} \emph{rasoḍũ}
  \item \emph{Say it:} \emph{rasoḍũ}, once more
  \item \emph{Say it:} \emph{mātā rasoḍāmā̃ rā̃dhe che} — Mother cooks in the kitchen.
  \item \emph{Recall:} say the Gujarati for to change, then the Gujarati for room, then say \emph{rasoḍũ} again
\end{itemize}

\begin{tcolorbox}[breakable,colback=teal!4,colframe=teal!35!black,title={Before you move on}]
What does \gu{રસોડું} mean? (\textquotedblleft{}Kitchen\textquotedblright{}.) Say it once more.
\end{tcolorbox}
\section[\texorpdfstring{\gu{અગાસી} (agāsī)}{અગાસી (agāsī)}]{\gu{અગાસી} (agāsī) — terrace, flat roof}

\label{lesson:GU-C153-agasi}



\begin{quote}
Before the new one: say the Gujarati for room, then the Gujarati for kitchen.
\end{quote}

\subsection*{You'll want to know: \gu{અગાસી}}
\textbf{\gu{અગાસી}} — \emph{agāsī} — \textquotedblleft{}terrace, flat roof\textquotedblright{}.

\begin{quote}
\textbf{\gu{ઉનાળામાં} \gu{અમે} \gu{અગાસીમાં} \gu{સૂઈએ} \gu{છીએ}.} — \emph{unāḷāmā̃ ame agāsīmā̃ sūīe chīe} — \textquotedblleft{}In summer we sleep on the terrace.\textquotedblright{}
\end{quote}

\begin{cousinweb}
The origin of \textbf{\gu{અગાસી}} (\emph{agāsī}) is not settled. CDIAL lists no Gujarati form of it, and EtymDB has no entry for it.
\end{cousinweb}

\begin{grammarlens}[title={What you've built}]
One more word for everyday life at A2.
\end{grammarlens}

\subsection*{Your turn}
\begin{itemize}
\raggedright
  \item \emph{Say it:} \emph{agāsī}
  \item \emph{Say it:} \emph{agāsī}, once more
  \item \emph{Say it:} \emph{unāḷāmā̃ ame agāsīmā̃ sūīe chīe} — In summer we sleep on the terrace.
  \item \emph{Recall:} say the Gujarati for room, then the Gujarati for kitchen, then say \emph{agāsī} again
\end{itemize}

\begin{tcolorbox}[breakable,colback=teal!4,colframe=teal!35!black,title={Before you move on}]
What does \gu{અગાસી} mean? (\textquotedblleft{}Terrace, flat roof\textquotedblright{}.) Say it once more.
\end{tcolorbox}
\section[\texorpdfstring{\gu{ઓસરી} (osrī)}{ઓસરી (osrī)}]{\gu{ઓસરી} (osrī) — veranda}

\label{lesson:GU-C153-osri}



\begin{quote}
Before the new one: say the Gujarati for kitchen, then the Gujarati for terrace.
\end{quote}

\subsection*{You'll want to know: \gu{ઓસરી}}
\textbf{\gu{ઓસરી}} — \emph{osrī} — \textquotedblleft{}veranda\textquotedblright{}.

\begin{quote}
\textbf{\gu{દાદા} \gu{ઓસરીમાં} \gu{બેસે} \gu{છે}.} — \emph{dādā osrīmā̃ bese che} — \textquotedblleft{}Grandfather sits on the veranda.\textquotedblright{}
\end{quote}

\begin{cousinweb}
\textbf{\gu{ઓસરી}} (\emph{osrī}) is inherited, but its source is not settled. CDIAL glosses it \textquotedblleft{}front portion of a house\textquotedblright{} and lists it twice: under \emph{*avaśara}, \textquotedblleft{}protection, shelter\textquotedblright{}, and under Sanskrit \emph{upasara}, \textquotedblleft{}approach\textquotedblright{} (CDIAL 848, 2258).
\end{cousinweb}

\begin{grammarlens}[title={What you've built}]
One more word for everyday life at A2.
\end{grammarlens}

\subsection*{Your turn}
\begin{itemize}
\raggedright
  \item \emph{Say it:} \emph{osrī}
  \item \emph{Say it:} \emph{osrī}, once more
  \item \emph{Say it:} \emph{dādā osrīmā̃ bese che} — Grandfather sits on the veranda.
  \item \emph{Recall:} say the Gujarati for kitchen, then the Gujarati for terrace, then say \emph{osrī} again
\end{itemize}

\begin{tcolorbox}[breakable,colback=teal!4,colframe=teal!35!black,title={Before you move on}]
What does \gu{ઓસરી} mean? (\textquotedblleft{}Veranda\textquotedblright{}.) Say it once more.
\end{tcolorbox}
\section[\texorpdfstring{\gu{દાદર} (dādar)}{દાદર (dādar)}]{\gu{દાદર} (dādar) — stairs, staircase}

\label{lesson:GU-C153-dadar}



\begin{quote}
Before the new one: say the Gujarati for terrace, then the Gujarati for veranda.
\end{quote}

\subsection*{You'll want to know: \gu{દાદર}}
\textbf{\gu{દાદર}} — \emph{dādar} — \textquotedblleft{}stairs, staircase\textquotedblright{}.

\begin{quote}
\textbf{\gu{તે} \gu{દાદર} \gu{ચઢે} \gu{છે}.} — \emph{te dādar caḍhe che} — \textquotedblleft{}He climbs the stairs.\textquotedblright{}
\end{quote}

\begin{cousinweb}
\textbf{\gu{દાદર}} (\emph{dādar}) is counted as inherited. CDIAL lists Gujarati \emph{dādar}, \textquotedblleft{}ladder\textquotedblright{}, under Sanskrit \emph{dardara}, \textquotedblleft{}split, broken\textquotedblright{}, without saying how the sense moved (CDIAL 6196).
\end{cousinweb}

\begin{grammarlens}[title={What you've built}]
One more word for everyday life at A2.
\end{grammarlens}

\subsection*{Your turn}
\begin{itemize}
\raggedright
  \item \emph{Say it:} \emph{dādar}
  \item \emph{Say it:} \emph{dādar}, once more
  \item \emph{Say it:} \emph{te dādar caḍhe che} — He climbs the stairs.
  \item \emph{Recall:} say the Gujarati for terrace, then the Gujarati for veranda, then say \emph{dādar} again
\end{itemize}

\begin{tcolorbox}[breakable,colback=teal!4,colframe=teal!35!black,title={Before you move on}]
What does \gu{દાદર} mean? (\textquotedblleft{}Stairs, staircase\textquotedblright{}.) Say it once more.
\end{tcolorbox}
